% Part: normal-modal-logic
% Chapter: axioms-systems

\documentclass[../../../include/open-logic-chapter]{subfiles}

\begin{document}

\olchapter{nml}{prf}{公理化推导}

\begin{explain}
在公理推导中，证明是一个公式序列。
序列中的每个公式要么是公理，要么是由序列中先前的公式经推理规则推导得出。
序列的最后一个公式即为证明的结论。
与自然演绎或矢列演算不同，这里不存在之后会被解除的“临时”假设；
任何所需的内容都必须表述为公理，或运用推理规则从公理推导得出。
\end{explain}

\olimport{introduction}
\olimport{normal-logics}
\olimport{logics-proofs}
\olimport{proofs-in-K}
\olimport{derived-rules}
\olimport{more-proofs-in-K}
\olimport{duals}
\olimport{proofs-modal-systems}
\olimport{soundness}
\olimport{systems-distinct}
\olimport{provability-from-set}
\olimport{provability-properties}
\olimport{consistency}

\OLEndChapterHook

\end{document}
